\section{Nueva frontera II: modelos de lenguaje grande (LLMs)}

\subsection{Definición y ubicación de los LLM}

\begin{figure}[H]
\centering
\includegraphics[width=0.85\textwidth]{slides-images/slide-69.jpg}
\caption{Ubicación de los LLM en el campo de la IA\protect\footnotemark}
\end{figure}
\footnotetext{Diapositiva 69.}

\begin{importantbox}{¿Qué es un modelo de lenguaje grande (LLM)?}
Los LLM son un concepto jerárquico dentro del campo de la IA:
\begin{itemize}
    \item \textbf{Artificial Intelligence}: todas las técnicas que hacen que una computadora simule el comportamiento humano
    \item \textbf{Deep Learning}: el uso de redes neuronales para extraer patrones de los datos
    \item \textbf{Large Language Models}: \textbf{redes neuronales de escala muy grande} entrenadas sobre \textbf{datos de texto de escala muy grande}
\end{itemize}
Palabras clave: ``very, very large neural networks trained on very, very large sets of text''.
\end{importantbox}

\subsection{Cómo funcionan los LLM: Next Token Prediction}

\begin{figure}[H]
\centering
\includegraphics[width=0.85\textwidth]{slides-images/slide-70.jpg}
\caption{Flujo de entrenamiento de un LLM tipo GPT\protect\footnotemark}
\end{figure}
\footnotetext{Diapositiva 70.}

El entrenamiento de un LLM puede resumirse en un objetivo sumamente sencillo:

$$
\text{Given a sequence of tokens, predict the next token.}
$$

Flujo concreto:
\begin{enumerate}
    \item \textbf{Conjunto de datos}: corpus de texto a gran escala como Common Crawl o WebText, divididos en tokens
    \item \textbf{Modelo}: por ejemplo GPT-3, con 175,000 millones de parámetros
    \item \textbf{Objetivo de entrenamiento}: dada una secuencia de tokens, predecir la distribución de probabilidad del siguiente token
    \item \textbf{Función de pérdida}: Cross-Entropy Loss, que mide la diferencia entre la distribución predicha y el siguiente token real
\end{enumerate}

\begin{figure}[H]
\centering
\includegraphics[width=0.85\textwidth]{slides-images/slide-71.jpg}
\caption{Flujo detallado de entrenamiento de Next Token Prediction\protect\footnotemark}
\end{figure}
\footnotetext{Diapositiva 71.}

\begin{knowledgebox}{De Next Token Prediction a la generación de texto}
Una vez concluido el entrenamiento, el LLM se usa mediante \textbf{generación autorregresiva} (autoregressive generation):
\begin{enumerate}
    \item El usuario introduce un prompt (por ejemplo, ``I'm giving a talk on AI at MIT. Can you outline it?'')
    \item El modelo predice la distribución de probabilidad del siguiente token y muestrea un token
    \item El nuevo token se agrega al final de la secuencia y se repite el proceso anterior
    \item Hasta generar un token de fin o alcanzar la longitud máxima
\end{enumerate}
Toda la ``inteligencia'' emergente se sustenta, en esencia, en esta tarea aparentemente simple de next token prediction.
\end{knowledgebox}

\subsection{Capacidades de los LLM}

\begin{figure}[H]
\centering
\includegraphics[width=0.85\textwidth]{slides-images/slide-73.jpg}
\caption{Capacidades actualmente confiables de los LLM\protect\footnotemark}
\end{figure}
\footnotetext{Diapositiva 73.}

GPT y otros LLM ya muestran un ``dominio'' (mastery over natural language) del lenguaje natural, con un desempeño confiable en los siguientes aspectos:
\begin{itemize}
    \item \textbf{Recuperación de conocimiento} (Knowledge Retrieval): responder preguntas factuales
    \item \textbf{Asistencia en escritura} (Writing Co-Pilot): redacción, pulido y traducción de textos
    \item \textbf{Asistencia en planeación} (Planning Co-Pilot): elaboración de planes y diseño de experimentos
\end{itemize}

\subsection{Limitaciones de los LLM}

\begin{figure}[H]
\centering
\includegraphics[width=0.85\textwidth]{slides-images/slide-74.jpg}
\caption{Las cuatro grandes limitaciones de los LLM\protect\footnotemark}
\end{figure}
\footnotetext{Diapositiva 74.}

\begin{warningbox}{Los cuatro retos clave de los LLM}
\begin{enumerate}
    \item \textbf{Robustez} (Robustness): son extremadamente sensibles a errores de ortografía o cambios de formato en la entrada (por ejemplo, ``Cn @uN66rN you translate ths from Spanish to English?'')
    \item \textbf{Alucinaciones} (Hallucinations): generan información incorrecta con total seguridad. ``GPT is a language model that...'' puede parecer perfectamente razonable, pero el contenido puede ser completamente inventado
    \item \textbf{Barandales y jailbreaks} (Guardrails and Jailbreaks): las restricciones de seguridad a nivel del modelo pueden eludirse con estrategias ingeniosas
    \item \textbf{Lógica y razonamiento numérico} (Logic and Numerics): tienen un desempeño débil en tareas que requieren razonamiento lógico riguroso y cálculo matemático
\end{enumerate}
\end{warningbox}

La raíz común de estos retos está en el \textbf{proceso de pensamiento de alto nivel} de los LLM: la robustez y la calibración de la confianza, la capacidad de planeación de largo plazo, y el razonamiento lógico y el descubrimiento científico son, todos ellos, fronteras actuales de investigación.

\subsection{Capacidades emergentes y leyes de escalamiento}

\begin{figure}[H]
\centering
\includegraphics[width=0.85\textwidth]{slides-images/slide-75.jpg}
\caption{Capacidades que emergen al crecer la escala del modelo\protect\footnotemark}
\end{figure}
\footnotetext{Diapositiva 75. Cita a Wei y otros, TMLR 2022.}

\begin{importantbox}{Capacidades emergentes (Emergent Abilities)}
Se dice que una capacidad es \textbf{emergente} (emergent) si no existe en modelos pequeños pero aparece en modelos grandes. En concreto:
\begin{itemize}
    \item \textbf{Estructuración del lenguaje} (Structuring Language): surge de manera repentina cuando el modelo alcanza del orden de $\sim 10^{22}$ parámetros
    \item \textbf{Comprensión de la fonética} (Understanding Phonetics): un salto igualmente abrupto
    \item \textbf{Realización de operaciones aritméticas} (Performing Arithmetic): aparece solo a mayor escala
\end{itemize}
La aparición de estas capacidades presenta un rasgo de ``transición de fase'' (phase transition): no es una mejora gradual, sino que surge de golpe al cruzar cierto umbral de escala.
\end{importantbox}

\begin{figure}[H]
\centering
\includegraphics[width=0.85\textwidth]{slides-images/slide-76.jpg}
\caption{Árbol de capacidades emergentes: a partir de 8,000 millones de parámetros\protect\footnotemark}
\end{figure}
\footnotetext{Diapositiva 76. Cita a Google AI Blog.}

A medida que la escala del modelo crece a partir de 8,000 millones de parámetros, el LLM va adquiriendo, en orden, comprensión del lenguaje, operaciones aritméticas y capacidad de responder preguntas, como un árbol de capacidades que crece y florece sin cesar.

\subsection{Resumen de la sección}

Los LLM son, en esencia, redes neuronales de escala muy grande entrenadas sobre volúmenes masivos de datos de texto, que logran una capacidad lingüística sorprendente mediante el objetivo simple de next token prediction. Sus capacidades presentan un carácter emergente conforme crece la escala del modelo. Pero los LLM aún enfrentan retos importantes en robustez, alucinaciones, protecciones de seguridad y razonamiento lógico, que son también las direcciones de investigación más activas en la actualidad.

%% =====================================================
